% =====================================================================
% anexo.tex — Anexo A: Repositório de Código-Fonte
% =====================================================================

\chapter{Repositório de Código-Fonte}
\label{an:repositorio}

O código-fonte completo do sistema descrito neste trabalho — incluindo
os módulos \texttt{main.py}, \texttt{kinect\_sensor.py},
\texttt{motor\_caixao\_areia.py} e \texttt{mde\_cartografia.py}, a
suíte de 83 testes unitários automatizados
(\texttt{test\_motor\_caixao.py}), a documentação técnica complementar
(\texttt{DOCUMENTACAO\_OFICIAL.md}, \texttt{GUIA\_CALIBRACAO\_PROJETOR.md},
\texttt{GUIA\_INSTALACAO\_INICIANTES.md}) e o
histórico completo de versões — está publicamente disponível no
seguinte repositório Git:

\begin{center}
\url{https://github.com/raquelbelchior1/PFC-2026}
\end{center}

\noindent O repositório contém, ainda, a pasta \texttt{installer/} com
o instalador \textit{offline} para Windows descrito na
Seção~\ref{subsec:instalador}, o manual do operador e os arquivos-fonte
deste próprio relatório em \LaTeX{} (pasta \texttt{relatorio\_tcc/}),
permitindo a reprodução integral do ambiente de desenvolvimento e da
suíte de testes a partir do código publicado.
